\chapter{继承、多态和类型擦除}

\section{问题入口：多个数据源如何统一使用}

假设一个服务需要读取用户信息。测试时从内存读取，生产环境从文件读取，未来可能从数据库读取。最差的写法是在业务逻辑里写分支：

\begin{lstlisting}[language=C++,caption={业务逻辑直接判断数据源类型}]
if (mode == "memory") {
    // 从内存找用户
} else if (mode == "file") {
    // 从文件找用户
} else if (mode == "database") {
    // 从数据库找用户
}
\end{lstlisting}

这个写法让业务逻辑知道太多环境细节。每新增一种数据源，都要修改业务代码。更好的方向是定义共同接口，让业务只依赖“能按 ID 查用户”这个能力。

\section{虚函数接口}

经典运行期多态用抽象基类：

\begin{lstlisting}[language=C++,caption={用户仓库接口}]
class UserRepository {
public:
    virtual ~UserRepository() = default;
    [[nodiscard]] virtual std::optional<UserRecord> find_by_id(int id) const = 0;
};
\end{lstlisting}

内存实现：

\begin{lstlisting}[language=C++,caption={内存仓库实现}]
class InMemoryUserRepository final : public UserRepository {
public:
    void add(UserRecord user)
    {
        this->users_[user.id] = std::move(user);
    }

    [[nodiscard]] std::optional<UserRecord> find_by_id(int id) const override
    {
        const auto found = this->users_.find(id);
        if (found == this->users_.end()) {
            return std::nullopt;
        }
        return found->second;
    }

private:
    std::unordered_map<int, UserRecord> users_;
};
\end{lstlisting}

业务服务只依赖接口：

\begin{lstlisting}[language=C++,caption={业务逻辑不关心数据源}]
class GreetingService {
public:
    explicit GreetingService(const UserRepository& repository)
        : repository_{repository}
    {
    }

    [[nodiscard]] std::string greeting_for(int id) const
    {
        const std::optional<UserRecord> user = this->repository_.find_by_id(id);
        if (!user.has_value()) {
            return "unknown user";
        }
        return "hello, " + user->name;
    }

private:
    const UserRepository& repository_;
};
\end{lstlisting}

虚函数的价值是运行期替换实现。测试和生产可以传不同对象，业务代码不用重新编译成模板实例。

\section{继承不是为了复用字段}

很多继承层次出问题，是因为把继承当成“共享成员变量”的工具。比如 \code{Cat} 和 \code{Dog} 都有名字，就让它们继承 \code{AnimalBase}。如果没有真正的多态行为，这种继承往往只会增加耦合。

优先组合：

\begin{lstlisting}[language=C++,caption={组合共享数据}]
struct NamedEntity {
    std::string name;
};

class Cat {
public:
    explicit Cat(std::string name)
        : identity_{std::move(name)}
    {
    }

private:
    NamedEntity identity_;
};
\end{lstlisting}

继承应当表达“可以通过基类接口使用派生对象”。如果没有这个需求，不要为了少写一个字段引入继承。

\section{虚析构函数为什么必要}

如果通过基类指针删除派生对象，基类析构函数必须是虚的：

\begin{lstlisting}[language=C++,caption={多态基类需要虚析构}]
class Plugin {
public:
    virtual ~Plugin() = default;
    virtual void run() = 0;
};
\end{lstlisting}

否则删除 \code{Plugin*} 时，派生类析构可能不执行，资源泄漏。只要一个类有虚函数并打算作为基类使用，虚析构通常就是必要的。

\section{override 和 final}

\code{override} 让编译器检查你确实重写了基类虚函数：

\begin{lstlisting}[language=C++,caption={override 防止签名写错}]
class FileUserRepository final : public UserRepository {
public:
    [[nodiscard]] std::optional<UserRecord> find_by_id(int id) const override;
};
\end{lstlisting}

如果你不小心少写 \code{const}，编译器会报错，而不是悄悄定义一个新函数。\code{final} 表示这个类不再作为基类继续派生，能让设计意图更明确。

\section{类型擦除：不暴露具体类型}

\code{std::function} 是一种类型擦除：调用者只知道它能按某个签名调用，不知道里面是函数指针、lambda 还是函数对象。我们也可以用它替代简单接口：

\begin{lstlisting}[language=C++,caption={用函数包装查找能力}]
using FindUser = std::function<std::optional<UserRecord>(int)>;

class GreetingService2 {
public:
    explicit GreetingService2(FindUser find_user)
        : find_user_{std::move(find_user)}
    {
    }

    [[nodiscard]] std::string greeting_for(int id) const
    {
        const std::optional<UserRecord> user = this->find_user_(id);
        if (!user.has_value()) {
            return "unknown user";
        }
        return "hello, " + user->name;
    }

private:
    FindUser find_user_;
};
\end{lstlisting}

这种方式适合只有一两个操作的小接口。操作很多、需要共享状态和清晰生命周期时，抽象基类可能更合适。类型擦除不是继承的替代宗教，而是另一个工具。

\section{variant：封闭集合的多态}

如果可能类型是封闭的，比如输入可以是文件、标准输入或内存字符串，可以用 \code{std::variant}：

\begin{lstlisting}[language=C++,caption={封闭类型集合用 variant}]
struct FileInput {
    std::filesystem::path path;
};

struct StdinInput {
};

struct MemoryInput {
    std::string text;
};

using InputSource = std::variant<FileInput, StdinInput, MemoryInput>;
\end{lstlisting}

处理时用 \code{std::visit}：

\begin{lstlisting}[language=C++,caption={访问 variant}]
std::string load_input(const InputSource& source)
{
    return std::visit([](const auto& item) -> std::string {
        return load_from(item);
    }, source);
}
\end{lstlisting}

新增一种输入类型时，编译器会帮助你找到需要处理的地方。虚函数适合开放集合，\code{variant} 适合封闭集合。先判断变化方向，再选机制。

\section{多态选择表}

\begin{longtable}{p{0.28\linewidth}p{0.28\linewidth}p{0.34\linewidth}}
\toprule
场景 & 机制 & 取舍 \\
\midrule
运行期替换多个实现 & 虚函数接口 & 稳定边界，有间接调用成本 \\
小型回调或策略 & \code{std::function} & 简单灵活，可能类型擦除开销 \\
编译期已知策略 & 模板参数 & 可内联，编译膨胀和报错更重 \\
封闭类型集合 & \code{std::variant} & 编译期穷尽检查，新增类型需改访问点 \\
共享少量字段 & 组合 & 耦合低，不提供运行期多态 \\
\bottomrule
\end{longtable}

\begin{exercisebox}
把问候服务分别用虚函数接口、\code{std::function} 和 \code{std::variant} 写三版。比较测试代码、生产组装代码和新增一种数据源时要改哪些地方。不要只看代码行数，要看依赖方向和错误暴露位置。
\end{exercisebox}
